\documentclass[english]{FUpowerdot}
\pagestyle{empty}
% Keep the starter neutral; users can add the FU Berlin marks with
% \framelogo and \titlelogo when their use is authorised.
\framelogo{}
\titlelogo{}

% Powerdot uses a DVI/PostScript workflow.  The class and style files in this
% starter are the complete FU Berlin presentation layer; the title and body
% below are intentionally neutral so they can be replaced without editing the
% institutional artwork.
\title{Research update}
\subtitle{A compact presentation starter}
\author{Alex Morgan}
\institute{Example Research Group}
\date{\today}
\fachbereich{Research}
\pdsetup{trans=none}

\begin{document}
\maketitle

\begin{slide}[toc=]{Overview}
  \tableofcontents
\end{slide}

\section{The question}
\begin{slide}{The question}
  \begin{itemize}
    \item State the decision or research question in one sentence.
    \item Give the audience the context needed to evaluate the answer.
  \end{itemize}
\end{slide}

\section{Evidence}
\begin{slide}{Evidence and next steps}
  \begin{itemize}
    \item Summarise the strongest evidence or result.
    \item Close with an action, open question, or discussion prompt.
  \end{itemize}
\end{slide}
\end{document}
